\scanpage{9}
% archive foliation: 8
\uline{Cor 1} \deleted{Dans} \deleted{\unsure{de plus}} ces conditions, on obtient encore
des conditions équivalentes en
exigeant \inserted{de plus} dans a) que les $L_n$, \deleted{$\pi_n(L_*)$}
soient de type fini, dans a') que les
$L_n^!$, \deleted{$H_n(L^!_*)$} soient de type fini, dans
b) et b') que les $M_i$ et $N_j$ soient de
type fini.

\unsure{On aura} sous les conditions de prop b
\[
L_n=\prod_{1\leq i\leq n+1}\unsure{\Lambda^i\Phi_{nk}}\otimes_k M_i
\times\prod_{\deleted{1}\inserted{0}\leq j\leq n}\unsure{\Lambda^j\Psi_{nk}}\otimes N_j
\]
\deleted{\unsure{donc on trouve}}

\uline{Cor 2} Sous les conditions équivalentes de
la proposition, pour que $L^!_*$ soit à degrés
\deleted{\unsure{finis par}} degrés bornés, il faut,
qu'il n'y ait qu'un nombre \inserted{fini} d'indices
$i$ et $j$ tels que $M_i\neq0$, $N_j\neq0$.

\deleted{Sous les conditions}
On appellera un $L_*$ \deleted{\unsure{spécial} \illegible{2}} \unsure{si}
$L^!_*$ est \deleted{\unsure{parfait}}, \inserted{à degrés bornés et à composantes}
et à $H_n$ projectifs de \unsure{t.f.|r.f.} --- i.e.
si $L_*$ est $\simeq$ somme finie de
modules de la forme \deleted{\unsure{$L$}}
\[
\Lambda^i\Phi_{\unsure{*}}\otimes_k M_i\inserted{\ (i\geq1)},\quad
\Lambda^j\Psi_{\unsure{*}}\otimes N_j\quad(j\geq0),
\]
avec les $M_i$ et $N_j$ proj. de \unsure{t.f.|r.f.}; \uline{stricte-}
\uline{ment \unsure{spéciaux|spécial}} si de plus les
$M_i$, $N_j$ sont \uline{libres}. Je reviens au cas des
